% Options for packages loaded elsewhere
\PassOptionsToPackage{unicode}{hyperref}
\PassOptionsToPackage{hyphens}{url}
\documentclass[
]{article}
\usepackage{xcolor}
\usepackage[margin=1in]{geometry}
\usepackage{amsmath,amssymb}
\setcounter{secnumdepth}{-\maxdimen} % remove section numbering
\usepackage{iftex}
\ifPDFTeX
  \usepackage[T1]{fontenc}
  \usepackage[utf8]{inputenc}
  \usepackage{textcomp} % provide euro and other symbols
\else % if luatex or xetex
  \usepackage{unicode-math} % this also loads fontspec
  \defaultfontfeatures{Scale=MatchLowercase}
  \defaultfontfeatures[\rmfamily]{Ligatures=TeX,Scale=1}
\fi
\usepackage{lmodern}
\ifPDFTeX\else
  % xetex/luatex font selection
\fi
% Use upquote if available, for straight quotes in verbatim environments
\IfFileExists{upquote.sty}{\usepackage{upquote}}{}
\IfFileExists{microtype.sty}{% use microtype if available
  \usepackage[]{microtype}
  \UseMicrotypeSet[protrusion]{basicmath} % disable protrusion for tt fonts
}{}
\makeatletter
\@ifundefined{KOMAClassName}{% if non-KOMA class
  \IfFileExists{parskip.sty}{%
    \usepackage{parskip}
  }{% else
    \setlength{\parindent}{0pt}
    \setlength{\parskip}{6pt plus 2pt minus 1pt}}
}{% if KOMA class
  \KOMAoptions{parskip=half}}
\makeatother
\usepackage{color}
\usepackage{fancyvrb}
\newcommand{\VerbBar}{|}
\newcommand{\VERB}{\Verb[commandchars=\\\{\}]}
\DefineVerbatimEnvironment{Highlighting}{Verbatim}{commandchars=\\\{\}}
% Add ',fontsize=\small' for more characters per line
\usepackage{framed}
\definecolor{shadecolor}{RGB}{248,248,248}
\newenvironment{Shaded}{\begin{snugshade}}{\end{snugshade}}
\newcommand{\AlertTok}[1]{\textcolor[rgb]{0.94,0.16,0.16}{#1}}
\newcommand{\AnnotationTok}[1]{\textcolor[rgb]{0.56,0.35,0.01}{\textbf{\textit{#1}}}}
\newcommand{\AttributeTok}[1]{\textcolor[rgb]{0.13,0.29,0.53}{#1}}
\newcommand{\BaseNTok}[1]{\textcolor[rgb]{0.00,0.00,0.81}{#1}}
\newcommand{\BuiltInTok}[1]{#1}
\newcommand{\CharTok}[1]{\textcolor[rgb]{0.31,0.60,0.02}{#1}}
\newcommand{\CommentTok}[1]{\textcolor[rgb]{0.56,0.35,0.01}{\textit{#1}}}
\newcommand{\CommentVarTok}[1]{\textcolor[rgb]{0.56,0.35,0.01}{\textbf{\textit{#1}}}}
\newcommand{\ConstantTok}[1]{\textcolor[rgb]{0.56,0.35,0.01}{#1}}
\newcommand{\ControlFlowTok}[1]{\textcolor[rgb]{0.13,0.29,0.53}{\textbf{#1}}}
\newcommand{\DataTypeTok}[1]{\textcolor[rgb]{0.13,0.29,0.53}{#1}}
\newcommand{\DecValTok}[1]{\textcolor[rgb]{0.00,0.00,0.81}{#1}}
\newcommand{\DocumentationTok}[1]{\textcolor[rgb]{0.56,0.35,0.01}{\textbf{\textit{#1}}}}
\newcommand{\ErrorTok}[1]{\textcolor[rgb]{0.64,0.00,0.00}{\textbf{#1}}}
\newcommand{\ExtensionTok}[1]{#1}
\newcommand{\FloatTok}[1]{\textcolor[rgb]{0.00,0.00,0.81}{#1}}
\newcommand{\FunctionTok}[1]{\textcolor[rgb]{0.13,0.29,0.53}{\textbf{#1}}}
\newcommand{\ImportTok}[1]{#1}
\newcommand{\InformationTok}[1]{\textcolor[rgb]{0.56,0.35,0.01}{\textbf{\textit{#1}}}}
\newcommand{\KeywordTok}[1]{\textcolor[rgb]{0.13,0.29,0.53}{\textbf{#1}}}
\newcommand{\NormalTok}[1]{#1}
\newcommand{\OperatorTok}[1]{\textcolor[rgb]{0.81,0.36,0.00}{\textbf{#1}}}
\newcommand{\OtherTok}[1]{\textcolor[rgb]{0.56,0.35,0.01}{#1}}
\newcommand{\PreprocessorTok}[1]{\textcolor[rgb]{0.56,0.35,0.01}{\textit{#1}}}
\newcommand{\RegionMarkerTok}[1]{#1}
\newcommand{\SpecialCharTok}[1]{\textcolor[rgb]{0.81,0.36,0.00}{\textbf{#1}}}
\newcommand{\SpecialStringTok}[1]{\textcolor[rgb]{0.31,0.60,0.02}{#1}}
\newcommand{\StringTok}[1]{\textcolor[rgb]{0.31,0.60,0.02}{#1}}
\newcommand{\VariableTok}[1]{\textcolor[rgb]{0.00,0.00,0.00}{#1}}
\newcommand{\VerbatimStringTok}[1]{\textcolor[rgb]{0.31,0.60,0.02}{#1}}
\newcommand{\WarningTok}[1]{\textcolor[rgb]{0.56,0.35,0.01}{\textbf{\textit{#1}}}}
\usepackage{longtable,booktabs,array}
\newcounter{none} % for unnumbered tables
\usepackage{calc} % for calculating minipage widths
% Correct order of tables after \paragraph or \subparagraph
\usepackage{etoolbox}
\makeatletter
\patchcmd\longtable{\par}{\if@noskipsec\mbox{}\fi\par}{}{}
\makeatother
% Allow footnotes in longtable head/foot
\IfFileExists{footnotehyper.sty}{\usepackage{footnotehyper}}{\usepackage{footnote}}
\makesavenoteenv{longtable}
\usepackage{graphicx}
\makeatletter
\newsavebox\pandoc@box
\newcommand*\pandocbounded[1]{% scales image to fit in text height/width
  \sbox\pandoc@box{#1}%
  \Gscale@div\@tempa{\textheight}{\dimexpr\ht\pandoc@box+\dp\pandoc@box\relax}%
  \Gscale@div\@tempb{\linewidth}{\wd\pandoc@box}%
  \ifdim\@tempb\p@<\@tempa\p@\let\@tempa\@tempb\fi% select the smaller of both
  \ifdim\@tempa\p@<\p@\scalebox{\@tempa}{\usebox\pandoc@box}%
  \else\usebox{\pandoc@box}%
  \fi%
}
% Set default figure placement to htbp
\def\fps@figure{htbp}
\makeatother
\setlength{\emergencystretch}{3em} % prevent overfull lines
\providecommand{\tightlist}{%
  \setlength{\itemsep}{0pt}\setlength{\parskip}{0pt}}
\usepackage{bookmark}
\IfFileExists{xurl.sty}{\usepackage{xurl}}{} % add URL line breaks if available
\urlstyle{same}
\hypersetup{
  pdftitle={Project\_Proposal},
  pdfauthor={Tim Eckart},
  hidelinks,
  pdfcreator={LaTeX via pandoc}}

\title{Project\_Proposal}
\author{Tim Eckart}
\date{2026-07-12}

\begin{document}
\maketitle

{
\setcounter{tocdepth}{2}
\tableofcontents
}
\subsection{Project Proposal}\label{project-proposal}

\emph{1.)} Load your dataset into R and take a snapshot of the data.
There is no requirement on how much to show, and there is no need to
show the entire dataset. Even 10 rows should be enough.

\begin{Shaded}
\begin{Highlighting}[]
\FunctionTok{library}\NormalTok{(data.table)}
\end{Highlighting}
\end{Shaded}

\begin{verbatim}
## 
## Attaching package: 'data.table'
\end{verbatim}

\begin{verbatim}
## The following object is masked from 'package:base':
## 
##     %notin%
\end{verbatim}

\begin{Shaded}
\begin{Highlighting}[]
\FunctionTok{library}\NormalTok{(readr)}
\NormalTok{azure\_vm\_raw }\OtherTok{=} \FunctionTok{read\_csv}\NormalTok{(}\StringTok{"trace\_data\_vmtable\_vmtable.csv"}\NormalTok{, }\AttributeTok{col\_names =} \ConstantTok{FALSE}\NormalTok{, }\AttributeTok{n\_max =} \DecValTok{10}\NormalTok{) }
\end{Highlighting}
\end{Shaded}

\begin{verbatim}
## Rows: 10 Columns: 11
\end{verbatim}

\begin{verbatim}
## -- Column specification --------------------------------------------------------
## Delimiter: ","
## chr (4): X1, X2, X3, X9
## dbl (7): X4, X5, X6, X7, X8, X10, X11
## 
## i Use `spec()` to retrieve the full column specification for this data.
## i Specify the column types or set `show_col_types = FALSE` to quiet this message.
\end{verbatim}

\begin{Shaded}
\begin{Highlighting}[]
\NormalTok{azure\_vm\_raw}
\end{Highlighting}
\end{Shaded}

\begin{verbatim}
## # A tibble: 10 x 11
##    X1        X2    X3        X4     X5      X6      X7      X8 X9      X10   X11
##    <chr>     <chr> <chr>  <dbl>  <dbl>   <dbl>   <dbl>   <dbl> <chr> <dbl> <dbl>
##  1 71fJw0x+~ GB6u~ 2sh/~ 5.58e5 1.67e6 91.8     0.729  20.8    Dela~     8    32
##  2 rKggHO/0~ ub4t~ +Zra~ 4.24e5 4.25e5 37.9     3.33   37.9    Unkn~     4    32
##  3 YrR8gPtB~ 9Lrd~ GEyI~ 1.13e6 1.13e6  0.304   0.221   0.304  Unkn~     4    32
##  4 xzQ++JF1~ 0XnZ~ 7aCQ~ 0      2.59e6 98.6    30.3    98.2    Inte~     2     4
##  5 vZEivnha~ HUGa~ /s/D~ 2.28e5 2.30e5 82.6    13.9    82.6    Unkn~     2     4
##  6 MqvcZ6Au~ p14c~ ZFCk~ 1.40e6 1.40e6  0.0979  0.0352  0.0979 Unkn~     4    32
##  7 034PavXA~ L9ut~ k2nh~ 1.42e6 1.42e6  0.0713  0.0327  0.0713 Unkn~     4    32
##  8 fBpt5HSY~ IwAB~ uYvK~ 2.41e6 2.41e6  0.260   0.0723  0.260  Unkn~    24    64
##  9 ZdSiRJc4~ 5tTf~ vLlC~ 1.66e5 1.68e5  0.0985  0.0342  0.0985 Unkn~     4    32
## 10 uzIY/rjn~ 3Gjj~ x9Ba~ 2.69e5 2.69e5  0.0662  0.0349  0.0662 Unkn~     4    32
\end{verbatim}

\begin{itemize}
\tightlist
\item
  \textbf{Answer:} Here is the data set listed from Microsoft Azure
  Public Dataset V2 on Github:
  \url{https://github.com/Azure/AzurePublicDataset/releases/tag/dataset-v2}
  . The data vmtable file has no header rows. The names are from the
  repository's schema file.
\end{itemize}

\begin{Shaded}
\begin{Highlighting}[]
\NormalTok{azure\_schema }\OtherTok{=} \FunctionTok{read\_csv}\NormalTok{(}\StringTok{"schema.csv"}\NormalTok{)}
\end{Highlighting}
\end{Shaded}

\begin{verbatim}
## Rows: 25 Columns: 4
## -- Column specification --------------------------------------------------------
## Delimiter: ","
## chr (3): filepattern, content, format
## dbl (1): field_number
## 
## i Use `spec()` to retrieve the full column specification for this data.
## i Specify the column types or set `show_col_types = FALSE` to quiet this message.
\end{verbatim}

\begin{Shaded}
\begin{Highlighting}[]
\NormalTok{knitr}\SpecialCharTok{::}\FunctionTok{kable}\NormalTok{(azure\_schema,}
             \AttributeTok{caption =} \StringTok{"Schema for the Azure 2019 VM trace (from schema.csv)"}\NormalTok{)}
\end{Highlighting}
\end{Shaded}

\begin{longtable}[]{@{}
  >{\raggedright\arraybackslash}p{(\linewidth - 6\tabcolsep) * \real{0.4953}}
  >{\raggedleft\arraybackslash}p{(\linewidth - 6\tabcolsep) * \real{0.1215}}
  >{\raggedright\arraybackslash}p{(\linewidth - 6\tabcolsep) * \real{0.2710}}
  >{\raggedright\arraybackslash}p{(\linewidth - 6\tabcolsep) * \real{0.1121}}@{}}
\caption{Schema for the Azure 2019 VM trace (from
schema.csv)}\tabularnewline
\toprule\noalign{}
\begin{minipage}[b]{\linewidth}\raggedright
filepattern
\end{minipage} & \begin{minipage}[b]{\linewidth}\raggedleft
field\_number
\end{minipage} & \begin{minipage}[b]{\linewidth}\raggedright
content
\end{minipage} & \begin{minipage}[b]{\linewidth}\raggedright
format
\end{minipage} \\
\midrule\noalign{}
\endfirsthead
\toprule\noalign{}
\begin{minipage}[b]{\linewidth}\raggedright
filepattern
\end{minipage} & \begin{minipage}[b]{\linewidth}\raggedleft
field\_number
\end{minipage} & \begin{minipage}[b]{\linewidth}\raggedright
content
\end{minipage} & \begin{minipage}[b]{\linewidth}\raggedright
format
\end{minipage} \\
\midrule\noalign{}
\endhead
\bottomrule\noalign{}
\endlastfoot
vm\_virtual\_core\_bucket\_definition.csv & 1 & bucket & STRING \\
vm\_virtual\_core\_bucket\_definition.csv & 2 & definition & STRING \\
vm\_memory\_bucket\_definition.csv & 1 & bucket & STRING \\
vm\_memory\_bucket\_definition.csv & 2 & definition & STRING \\
subscriptions/subscriptions.csv.gz & 1 & subscription id &
STRING\_HASH \\
subscriptions/subscriptions.csv.gz & 2 & timestamp first vm created &
INTEGER \\
subscriptions/subscriptions.csv.gz & 3 & count vms created & INTEGER \\
deployment/deployment.csv.gz & 1 & deployment id & STRING\_HASH \\
deployment/deployment.csv.gz & 2 & deployment size & INTEGER \\
vmtable/vmtable.csv.gz & 1 & vm id & STRING\_HASH \\
vmtable/vmtable.csv.gz & 2 & subscription id & STRING\_HASH \\
vmtable/vmtable.csv.gz & 3 & deployment id & STRING\_HASH \\
vmtable/vmtable.csv.gz & 4 & timestamp vm created & INTEGER \\
vmtable/vmtable.csv.gz & 5 & timestamp vm deleted & INTEGER \\
vmtable/vmtable.csv.gz & 6 & max cpu & DOUBLE \\
vmtable/vmtable.csv.gz & 7 & avg cpu & DOUBLE \\
vmtable/vmtable.csv.gz & 8 & p95 max cpu & DOUBLE \\
vmtable/vmtable.csv.gz & 9 & vm category & STRING \\
vmtable/vmtable.csv.gz & 10 & vm virtual core count bucket & STRING \\
vmtable/vmtable.csv.gz & 11 & vm memory (gb) bucket & STTRING \\
vm\_cpu\_readings/vm\_cpu\_readings-file-*-of-195.csv.gz & 1 & timestamp
& INTEGER \\
vm\_cpu\_readings/vm\_cpu\_readings-file-*-of-195.csv.gz & 2 & vm id &
STRING\_HASH \\
vm\_cpu\_readings/vm\_cpu\_readings-file-*-of-195.csv.gz & 3 & min cpu &
DOUBLE \\
vm\_cpu\_readings/vm\_cpu\_readings-file-*-of-195.csv.gz & 4 & max cpu &
DOUBLE \\
vm\_cpu\_readings/vm\_cpu\_readings-file-*-of-195.csv.gz & 5 & avg cpu &
DOUBLE \\
\end{longtable}

\begin{itemize}
\tightlist
\item
  Here is the Schema.csv listing the column headers name for the table.
\end{itemize}

\begin{Shaded}
\begin{Highlighting}[]
\FunctionTok{library}\NormalTok{(readr)}

\NormalTok{vm\_cols }\OtherTok{=} \FunctionTok{c}\NormalTok{(}\StringTok{"vm\_id"}\NormalTok{, }\StringTok{"subscription\_id"}\NormalTok{, }\StringTok{"deployment\_id"}\NormalTok{, }\StringTok{"timestamp\_vm\_created"}\NormalTok{, }\StringTok{"timestamp\_vm\_deleted"}\NormalTok{, }\StringTok{"max\_cpu"}\NormalTok{, }
            \StringTok{"avg\_cpu"}\NormalTok{, }\StringTok{"p95\_max\_cpu"}\NormalTok{, }\StringTok{"vm\_category"}\NormalTok{, }\StringTok{"vm\_virtual\_core\_count\_bucket"}\NormalTok{, }\StringTok{"vm\_memory\_gb\_bucket"}\NormalTok{)}

\NormalTok{azure\_vm\_table }\OtherTok{=} \FunctionTok{read\_csv}\NormalTok{(}\StringTok{"trace\_data\_vmtable\_vmtable.csv"}\NormalTok{, }\AttributeTok{col\_names =}\NormalTok{ vm\_cols, }
                          \AttributeTok{col\_types =} \FunctionTok{cols}\NormalTok{(}\AttributeTok{vm\_virtual\_core\_count\_bucket =} \FunctionTok{col\_character}\NormalTok{(), }\AttributeTok{vm\_memory\_gb\_bucket =} \FunctionTok{col\_character}\NormalTok{()))}

\FunctionTok{head}\NormalTok{(azure\_vm\_table[, }\SpecialCharTok{{-}}\FunctionTok{c}\NormalTok{(}\DecValTok{4}\SpecialCharTok{:}\DecValTok{11}\NormalTok{)],}\DecValTok{10}\NormalTok{)}
\end{Highlighting}
\end{Shaded}

\begin{verbatim}
## # A tibble: 10 x 3
##    vm_id                                           subscription_id deployment_id
##    <chr>                                           <chr>           <chr>        
##  1 71fJw0x+SDRdAxKPwLyHZhTgQpYw2afS6tjJhfT6kHnmLH~ GB6uQC1NSArW5n~ 2sh/ZjaYdfps~
##  2 rKggHO/04j31UFy65mDTwtjdMQL/G03xWfl3xGeiilB4/W~ ub4ty8ygwOECrI~ +ZraIDUNaWYD~
##  3 YrR8gPtBmfNaOdnNEW5If1SdTqQgGQHEnLHGPjySt53bKW~ 9LrdYRcUfGbmL2~ GEyIElfPSFup~
##  4 xzQ++JF1UAkh70CDhmzkiOo+DQn+E2TLErCFKEmSswv1pl~ 0XnZZ8sMN5HY+Y~ 7aCQS6fPUw9r~
##  5 vZEivnhabRmImDr+JqKqZnpIM3WxtypwoxjfjnklR/idyR~ HUGaZ+piPP4eHj~ /s/D5VtTQDxy~
##  6 MqvcZ6Au5ouI6if56MJHmoSqHtX8oRv0dPkaxCId3aUcr1~ p14cXGYqCKCcF7~ ZFCk80sIQzr4~
##  7 034PavXAkWso23QmI+NMty3u9XA2FVhIKKOZV+tbQru/n5~ L9utvn29uJD9za~ k2nh5l37CN78~
##  8 fBpt5HSYGsQyRQSSrhj+n5IDFzMJKf2k6Am1OFEIGpFl6n~ IwABYzfINELEB6~ uYvK2XjSSN+Q~
##  9 ZdSiRJc4t7SuV2uOfhbjkGXJPK6w1cuE8NcEYb+LJTOvWa~ 5tTf4IJxkp/u62~ vLlC4aSY52U7~
## 10 uzIY/rjnHYDAGJ7AHxffw+n+QlBriI1s12yb7MtYSykiZf~ 3GjjR481xxGvEc~ x9BaDDFxkejZ~
\end{verbatim}

\begin{Shaded}
\begin{Highlighting}[]
\FunctionTok{head}\NormalTok{(azure\_vm\_table[, }\SpecialCharTok{{-}}\FunctionTok{c}\NormalTok{(}\DecValTok{1}\SpecialCharTok{:}\DecValTok{3}\NormalTok{)],}\DecValTok{10}\NormalTok{)}
\end{Highlighting}
\end{Shaded}

\begin{verbatim}
## # A tibble: 10 x 8
##    timestamp_vm_created timestamp_vm_deleted max_cpu avg_cpu p95_max_cpu
##                   <dbl>                <dbl>   <dbl>   <dbl>       <dbl>
##  1               558300              1673700 91.8     0.729      20.8   
##  2               424500               425400 37.9     3.33       37.9   
##  3              1133100              1133700  0.304   0.221       0.304 
##  4                    0              2591400 98.6    30.3        98.2   
##  5               228300               229800 82.6    13.9        82.6   
##  6              1395600              1397700  0.0979  0.0352      0.0979
##  7              1422300              1422600  0.0713  0.0327      0.0713
##  8              2414400              2414700  0.260   0.0723      0.260 
##  9               165900               168300  0.0985  0.0342      0.0985
## 10               268500               269100  0.0662  0.0349      0.0662
## # i 3 more variables: vm_category <chr>, vm_virtual_core_count_bucket <chr>,
## #   vm_memory_gb_bucket <chr>
\end{verbatim}

\begin{Shaded}
\begin{Highlighting}[]
\FunctionTok{head}\NormalTok{(azure\_vm\_table[, }\SpecialCharTok{{-}}\FunctionTok{c}\NormalTok{(}\DecValTok{1}\SpecialCharTok{:}\DecValTok{8}\NormalTok{)],}\DecValTok{10}\NormalTok{)}
\end{Highlighting}
\end{Shaded}

\begin{verbatim}
## # A tibble: 10 x 3
##    vm_category       vm_virtual_core_count_bucket vm_memory_gb_bucket
##    <chr>             <chr>                        <chr>              
##  1 Delay-insensitive 8                            32                 
##  2 Unknown           4                            32                 
##  3 Unknown           4                            32                 
##  4 Interactive       2                            4                  
##  5 Unknown           2                            4                  
##  6 Unknown           4                            32                 
##  7 Unknown           4                            32                 
##  8 Unknown           24                           64                 
##  9 Unknown           4                            32                 
## 10 Unknown           4                            32
\end{verbatim}

\begin{itemize}
\tightlist
\item
  Here is the columns renamed to their respected headers per Schema.
\end{itemize}

\textbf{2.)} Write a paragraph saying who is working on the project (if
team), or your name otherwise. Explain briefly what the dataset is
about. Provide a link so I can access the dataset if necessary. State
how many observations there are.\n

Finally, write one or two sentences about what the goal of the analysis
is. Clearly state the response variable and explanatory variables. This
should be a simple thing such as: The goal of this analysis is to
predict the mpg for certain cars based on x,y,z,a,b,c,\ldots{}\n

\textbf{Answer:} This project will be conducted individually by Timothy
Eckart, using the \emph{vmtable} trace from the \textbf{Azure Public
Dataset V2}, published by Microsoft in the AzurePublicDataset repository
on GitHub
(\url{https://github.com/Azure/AzurePublicDataset/releases/tag/dataset-v2}).
The dataset is a sanitized, representative subset of Microsoft Azure's
first-party virtual machine workload in one geographical region,
collected over 30 consecutive days in 2019 and released under a CC-BY
license. The data contains 2,695,548 observations on virtual machines
from 6,687 Azure customer subscriptions, including average, maximum, and
95th-percentile CPU utilization, VM category, virtual core count, memory
size, and creation/deletion time stamps. Each observation represents one
virtual machine's full lifetime within the trace. The goal of this
analysis is to predict a virtual machine's average CPU utilization
(\texttt{avg\_cpu}) based on its category, virtual core count, memory
size, and lifetime derived from its creation and deletion timestamp. \n

\begin{itemize}
\tightlist
\item
  Per the dataset's CC-BY attribution requirement this project cites the
  accompanying SOSP'17 paper (Cortez 2017; see References).
\end{itemize}

\subsection{Reference}\label{reference}

Cortez, E., Bonde, A., Muzio, A., Russinovich, M., Fontoura, M., \&
Bianchini, R. (2017). Resource Central: Understanding and predicting
workloads for improved resource management in large cloud platforms. In
\emph{Proceedings of the 26th Symposium on Operating Systems Principles
(SOSP '17)} (pp.~153--167). ACM.
\url{https://www.microsoft.com/en-us/research/wp-content/uploads/2017/10/Resource-Central-SOSP17.pdf}
\n

Microsoft Azure. (2019). \emph{AzurePublicDataset: Microsoft Azure
traces} (Version 2) {[}Data set{]}. GitHub.
\url{https://github.com/Azure/AzurePublicDataset}

\n
\n

\end{document}
